\section{Fundamentação Teórica e Trabalhos Relacionados}

\subsection{Sistema de Templates MASTeR e Diretrizes INCOSE}
O sistema MASTeR~\cite{sophist2020master, rupp2014anforderungsschablonen} estrutura requisitos em linguagem controlada combinando um \textbf{template principal} obrigatório e um \textbf{sub-template de condição} opcional, vinculados a verbos modais estritos (\textit{shall}, \textit{should} ou \textit{will}). Os templates principais dividem-se em requisitos funcionais (\textit{FunctionMASTeR}, subdividido em atividade independente do sistema, interação com usuário e requisito de interface) e não funcionais (\textit{PropertyMASTeR}, \textit{EnvironmentMASTeR} e \textit{ProcessMASTeR}). Quando aplicável, a cláusula condicional precede o template principal segundo três padrões: \textit{LogicMASTeR} (iniciado por \textit{``If \dots''}), \textit{EventMASTeR} (\textit{``As soon as \dots''}) ou \textit{TimeMASTeR} (\textit{``As long as \dots''}), exigindo repetição explícita da palavra-chave condicional após operadores lógicos \textit{and}/\textit{or}.

Complementarmente à estrutura sintática, o \textit{INCOSE Guide to Writing Requirements}~\cite{incose2023guide} estabelece nove características individuais para requisitos bem formados (C1--C9): \textit{necessário}, \textit{apropriado}, \textit{não ambíguo}, \textit{completo}, \textit{singular}, \textit{viável}, \textit{verificável}, \textit{correto} e \textit{conforme}. Estudos empíricos de Gro\ss er et al.~\cite{grosser2023comparative, grosser2024benchmarking} confirmam que o uso do MASTeR eleva significativamente a singularidade (atomicidade) e a conformidade normativa frente ao texto livre, embora exija treinamento especializado para redação manual.

\subsection{Conversão Automática de Requisitos com LLMs}
Abordagens clássicas de reescrita automática de requisitos baseavam-se em regras gramaticais manuais, ontologias e árvores sintáticas~\cite{arora2015automated}. Com o advento dos LLMs, Norheim e Rebentisch~\cite{norheim2024structuring} exploraram preenchimento de lacunas (\textit{slot-filling}) via \textit{one-shot prompting} para o template EARS, porém sem avaliação sistemática em larga escala.

Ro\ss kothen et al.~\cite{rosskothen2026converting} avançaram o estado da arte ao propor a geração integral de sentenças MASTeR com decomposição livre ($1 \rightarrow n$ requisitos atômicos) controlada por decodificação estruturada (\textit{Constrained Decoding}) e refinada em um \textit{Reflection Loop}~\cite{ji2023mitigating}. Avaliando quatro configurações sobre 249 requisitos de cinco projetos reais reescritos por humanos em 410 sentenças MASTeR~\cite{grosser2023comparative}, os autores mostraram que métricas baseadas em \textit{BERTScore}~\cite{zhang2020bertscore}, \textit{Embedding Mover's Distance} (EMD, derivada de WMD~\cite{kusner2015word}) e conformidade de sub-templates apresentam correlação positiva moderada ($\rho = 0{,}503$ combinada) com o julgamento de especialistas da indústria.

\subsection{LLM-as-a-Judge Generativo vs. Modelos Discriminativos System One (JEV)}
No paradigma convencional de \textit{LLM-as-a-Judge}, um modelo generativo recebe instruções extensas e produz cadeias de raciocínio ou justificativas textuais seguidas de um veredito booleano ou numérico. Embora flexível, essa estratégia sofre de três gargalos em pipelines de Engenharia de Requisitos em lote: (i) \textbf{alta latência e custo de geração}, pois a decodificação autorregressiva é sequencial por token; (ii) \textbf{variância e deriva de justificativas}, exigindo estágios extras de filtragem para impedir que críticas excessivas induzam alucinações no modelo gerador; e (iii) \textbf{má calibração probabilística}, dificultando a definição de limiares confiáveis de aceitação automática.

Em contrapartida, modelos discriminativos de decisão denominados \textit{System One}, exemplificados pela arquitetura \textbf{JEV} (\textit{TypeSafe AI})~\cite{typesafe2026jev}, separam o julgamento estruturado da geração de prosa. Dado um contexto de entrada (\texttt{state}), o JEV avalia múltiplas consultas tipadas em um único passe paralelo (\textit{forward pass}), sem gerar tokens de texto livre:
\begin{itemize}
    \item \texttt{Choice}: classificação multiclasse mutuamente exclusiva retornando a classe escolhida, a distribuição de probabilidades sobre todas as opções e um índice de confiança calibrado;
    \item \texttt{Noul}: verificação binária de condição (\textit{checklist}), retornando diretamente a probabilidade calibrada $p \in [0, 1]$ de atendimento à regra (onde $p \approx 0{,}5$ indica alta incerteza epistêmica);
    \item \texttt{Score}: avaliação ordinal sobre níveis descritivos ordenados, retornando a posição contínua ponderada pelas probabilidades de cada nível e sua respectiva confiança.
\end{itemize}
Essa distinção arquitetural permite executar auditorias multi-critério em dezenas a poucas centenas de milissegundos, viabilizando políticas formais de cascata (\textit{Accept or Escalate}) em fluxos industriais de requisitos.
